%! Transformations of Functions — Unit 1, Lesson 1.4 — Warm-Up
\documentclass[10pt]{article}
\usepackage{precalculus-article}
\usepackage{precalculus-boxes}
\newcommand{\TallMath}[1]{$\displaystyle #1\rule[-1.4em]{0pt}{3.2em}$}

\begin{document}

\pageheader{Unit 1, Lesson 1.4}{Warm-Up}
\namedateperiod

\vspace{0.12in}

\begin{enumerate}[itemsep=14pt]

\item Let $f(x) = x^2 - 3$. Evaluate each expression.

      \medskip
      \begin{enumerate}[label=(\alph*), itemsep=6pt]
        \item $f(0) = $ \blank{3cm}
        \item $f(-2) = $ \blank{3cm}
        \item $f(x+1) = $ \blank{5cm}
      \end{enumerate}

\item The table below gives values of a function $f$.

      \medskip
      \renewcommand{\arraystretch}{1.3}
      \begin{tabular}{|c|c|c|c|c|c|}
      \hline
      $x$    & $-2$ & $-1$ & $0$ & $1$ & $2$ \\
      \hline
      $f(x)$ & $5$  & $2$  & $1$ & $2$ & $5$ \\
      \hline
      \end{tabular}

      \medskip
      State the domain and range of $f$ using the values in the table.

      \vspace{0.08in}
      Domain: \blank{5cm} \quad Range: \blank{5cm}

\item The graph of $p(x)$ is a polynomial with the following features: it passes through
      $(-2,\,0)$, $(1,\,0)$, and $(0,\,-2)$, and its end behavior is
      $p(x) \to +\infty$ as $x \to +\infty$ and $p(x) \to -\infty$ as $x \to -\infty$.

      Circle the best description of the degree and leading coefficient of $p(x)$.

      \medskip
      \textbf{A.}\ Even degree, positive leading coefficient \qquad
      \textbf{B.}\ Odd degree, positive leading coefficient

      \medskip
      \textbf{C.}\ Even degree, negative leading coefficient \qquad
      \textbf{D.}\ Odd degree, negative leading coefficient

\end{enumerate}

\end{document}
